\section{高迁移：把“收购前委员会”搬到公开市场}

个人能复制的不是买断合同，而是更严苛的入场条件。公开市场可能提供较低的进入门槛和可选择不交易的自由，但费用、税费、流动性、生活现金需求、账户条件和行为偏差会显著改变实际体验；它们不构成天然优势。个人没有控制权，信息不对称仍在，且行情波动会持续考验纪律。因此，最有价值的私募式能力是下列六项：

\begin{enumerate}
  \item \textbf{能力圈收窄：}只研究少量、能持续跟踪的产业链；具体数量应由可用时间、披露频率和理解深度决定，而非追逐所有热门主题。
  \item \textbf{现金流先于叙事：}利润、经营现金流、应收、维持性资本开支、债务和再融资共同决定下行。
  \item \textbf{运营改善必须可见：}把提价、产品结构、降本、产能利用率、去杠杆或资产盘活写成可追踪的关键经营指标（KPI），而不是愿景。
  \item \textbf{价值桥而非单点估值：}把收益拆成盈利/现金流兑现、分红或分派、以及可能的估值变化；后者不能被默认。
  \item \textbf{分期拨款：}证据增强才增加承诺；股价下跌不是增加承诺的理由。
  \item \textbf{反方与失效：}在买入前写下最可能推翻自身结论的事实、日期和行动，而不是在下跌后临时解释。
\end{enumerate}

\section{部分迁移：能观察，不能推动}

运营改善、实物资产和事件驱动都可以成为公开市场研究对象，但其本质发生了变化。私募投资者可以任命管理层、调整预算或谈判融资；个人只能判断管理层已经做了什么、是否正在兑现，以及价值是否已被价格充分反映。对此，正确的估值处理是保守的：\textbf{没有披露、没有量化桥、没有关键经营指标（KPI）的“改善”，估值信用为零。}

\begin{exhibitbox}[Exhibit 3：迁移阶梯]
\small
\begin{tabularx}{\exhibitboxwidth}{L{2.7cm} L{3.0cm} X L{3.1cm}}
\toprule
\textbf{类别} & \textbf{机构原型} & \textbf{个人可做的动作} & \textbf{不可跨越的边界} \\
\midrule
高度可迁移 & KKR / Apollo 式承保 & 现金流、债务、再融资、估值、KPI、反方审查。 & 不能把“会做尽调”误认为拥有专属信息。 \\
高度可迁移 & 私募投资委员会（IC）/ 资本调用 & 以公开证据分阶段复核研究结论。 & 不因价格更低而自动加仓。 \\
部分可迁移 & Bain / Blackstone 式运营改善 & 选择已公开启动的经营改变并跟踪结果。 & 无法强迫管理层执行，不能把计划当利润。 \\
部分可迁移 & Brookfield 式实物资产 & 用项目现金流、维护资本开支和融资压力研究公募 REIT。 & 没有项目控制权，也不应把分派当无风险利息。 \\
部分可迁移 & 正式公告事件研究 & 围绕公告、业绩、资产处置或去杠杆建立可为零的研究档案。 & 传闻、未公开信息和“消息灵通”不属于研究优势。 \\
不迁移 & 杠杆收购（LBO）、控制权、私贷、保险资产负债管理（ALM） & 无。 & 法律结构、资金成本、管理资源与流动性容忍度均不同。 \\
\bottomrule
\end{tabularx}
\sourcenote{S01--S07；制度转换为本报告分析。}
\end{exhibitbox}

\section{明确的红线}

\begin{riskbox}[这不是个人的“待学习清单”]
以下内容不应被包装为个人可照搬的方法：控制权收购与私有化、董事会影响、杠杆收购（LBO）融资、困境债权人控制、私募直贷发起、基金二级份额交易、需要稳定借券的多空策略、并购套利基础设施，以及任何非公开信息。它们依赖法律地位、资金池、运营团队、合规系统、交易对手与承受长期锁定的能力。缺少这些条件时，模仿只会放大尾部风险。
\end{riskbox}

\section{个人真正可能拥有的比较优势}

个人没有机构的访问权，却可能在个人现金需求、账户条件、交易成本和行为约束都允许时，选择不强制配置、以可为零的研究范围等待公开证据，并在真正理解的领域长期积累。这些可能性不是天然优势，无法保证超额收益，却有助于避免“看不懂也参与”。若投资者没有兴趣或时间维护这套流程，合理的结果可以是直接研究卫星为零，并在费用、账户、流动性和税费均核验后考虑更简单的分散化方案。
